% page 189 of the facsimile (image p-188 of the scan)
% block 175.3 x 237.7 mm ; left margin 26.7 mm
\pgc{5.080mm}{0.000mm}{
\l{\cel{58}181}
\l{}
\l{}
\l{fronto. I. La parto supra di la vizajo homala, inter la brovi e}
\l{\cel{3}la lineo ube komencas la hararo.- II. La parto supra di la vizajo}
\l{\cel{3}di la animali. - III. La parto supra di edifico, di monto, di}
\l{\cel{3}arboro. - IV. La facio avana di ula kozi (armeo qua marchas, ec.)}
\l{\cel{3}- DEFIRS.}
\l{}
\l{\sou{frontiero}.\cel{2}Limito qua separas la teritorio di stato de olta di}
\l{\cel{3}stato kontigua. - II. Limito qua impedas expanso. - EFIS.}
\l{}
\l{\sou{frontispico}. I. (\sou{arkitekt.)}\cel{2}Facio precipua di edifico. - II. Titulo}
\l{\cel{3}di libro, akompanata da vinyeti, e c. ; graburo lokizita opoze a}
\l{\cel{3}la pagino qua kontenas la titulo. - DEFIS.}
\l{}
\l{\sou{frontono}. Ornivo qua quaze kronizas la enireyo precipue di edifico,}
\l{\cel{3}e qua havas kom bazo la kornico di la entablamento.-EFIRS.}
\l{}
\l{\sou{frostar}. (\sou{netrans.}) Transformesar a glacio. - DE.}
\l{}
\l{\sou{frotar}.\cel{2}(\sou{trans}.)Igar (korpo) kontaktar kun altra korpo,quan onu}
\l{\cel{3}igas pasar sur ita presante. - DFS.}
\l{}
\l{\sou{frua}. Ante la instanto kande onu kustumas agar lo ordinare ; ante}
\l{\cel{3}la instanto kande onu expektas lua evento. - D. (+\sou{anticipa})}
\l{}
\l{\sou{frugala}. Qua saciesas per nutrivi simpla. - DEFIS.}
\l{}
\l{\sou{frugilego.(zool.}) Korvo di qua la beko esas rekta, akuta, ne garnisi-}
\l{\cel{3}ta ye plumi. - L.\cel{1}\sou{corvus frugilegus}.}
\l{}
\l{\sou{frukto}. (\sou{bot.)}\cel{2}I. Produkturo de la vejetanto, qua sucedas a floro.}
\l{\cel{3}II. Organo di la vejetanti (ovario, fekundigita) qua kreskas en}
\l{\cel{3}la floro, ube lu existas en estado rudimenta e kontenas la}
\l{\cel{3}semino (ovulo) destinita a reproduktar la planto. - II. (\sou{metaf}.)}
\l{\cel{3}a. (\sou{aludante patrino)}. La filio quan elu konceptis, quan elu}
\l{\cel{3}portas en sua ventro.\cel{2}b. (\sou{aludante mariajo}).La filio qua naskas}
\l{\cel{3}de la mariajo.\cel{2}c. Rezultajo avantajoza de ulo. - DEFIRS.}
\l{}
\l{\sou{frumento}. (\sou{bot.}) Planto cereala, de la familio \textquotedbl{}cereali\textquotedbl{}, di qua}
\l{\cel{3}la spiko esas kompozaja, kontenanta grano manjebla de qua onu}
\l{\cel{3}facas farino. - L.\cel{1}\sou{triticum}. - DEFIRS.}
\l{}
\l{\sou{frunsa}r. (\sou{trans.}) I. Plisar (la brovi) per kontrakto. - II. Igar}
\l{\cel{3}(stofo) plisita per pasar en lu filo o tre adheranta rubando.-FS.}
\l{}
\l{\sou{frusta}. Di qua la reliefo esas kruda, grosiera, vulgara. - FIS.}
\l{}
\l{\sou{frustrar}. (\sou{trans.,}\cel{1}de)\cel{2}Privacar (ulu) de havo, de avantajo quan}
\l{\cel{3}lu havas la yuro posedar. - FIS.}
\l{}
\l{ftizio. Deperiso pro morbo. - DEFIS.}
\l{}
\l{fudro. Tre granda barelo (de 800 litri e mem plu multe).- DF.}
}
